\documentclass[11pt,a4paper]{article}

% ── Packages ──
\usepackage[utf8]{inputenc}
\usepackage[T1]{fontenc}
\usepackage{amsmath,amssymb}
\usepackage{physics}
\usepackage{siunitx}
\usepackage{booktabs}
\usepackage{geometry}
\usepackage{hyperref}
\usepackage{enumitem}
\usepackage{float}
\usepackage{graphicx}

\geometry{margin=2.5cm}
\hypersetup{colorlinks=true,linkcolor=blue,citecolor=blue,urlcolor=blue}

\title{Residual Gas Velocity Above a Stopped 2-Arm Propeller\\
       in Dense Xenon: A Mechanism-by-Mechanism Analysis}
\author{}
\date{}

\begin{document}
\maketitle

% ══════════════════════════════════════════════════════════════════════════════
\section{Problem Definition}
\label{sec:problem}
% ══════════════════════════════════════════════════════════════════════════════

A horizontal 2-arm propeller (NACA\,0012 airfoil cross-section) rotates
inside a vertical cylindrical vessel filled with xenon gas at high pressure.
The propeller sweeps through $\theta = \pi$ (180°) via a bang-bang motion
profile. After rotation, the carrier plate slides radially and then lifts
vertically. We seek the residual gas velocity above the lifted plate to
determine when quiescence is achieved for ion detection.

\subsection{Parameters}

\begin{table}[H]
\centering
\caption{System parameters for the 2-arm propeller.}
\label{tab:params}
\begin{tabular}{llll}
\toprule
Symbol & Description & Value & Units \\
\midrule
$n$ & Number of blade arms & 2 & --- \\
$R$ & Blade span (hub to tip) & 1.6 & m \\
$c$ & Blade chord length & 0.16 & m \\
$\theta$ & Sweep angle ($= \pi$) & 180 & ° \\
$m_\mathrm{blade}$ & Mass of each blade & 0.25 & kg \\
$m_\mathrm{plate}$ & Mass of each carrier plate & 0.25 & kg \\
$\tau$ & Motor torque & 60 & N$\cdot$m \\
\midrule
\multicolumn{4}{l}{\textit{Carrier plate parameters}} \\
$r_0$ & Initial radial position & 0.8 & m \\
$F_\mathrm{act}$ & Linear actuator force & 20 & N \\
$L_c$ & Carrier width & 0.16 & m \\
\midrule
\multicolumn{4}{l}{\textit{Phase 3 parameters}} \\
$z_\mathrm{lift}$ & Lift height & 5 & mm \\
$t_\mathrm{lift}$ & Lift duration & 1.0 & s \\
$t_\mathrm{wait}$ & Wait time after lift & 0.5 & s \\
$u_c$ & Quiescence cutoff velocity & 0.1 & mm/s \\
\midrule
\multicolumn{4}{l}{\textit{Gas conditions}} \\
$P$ & Xenon pressure & 15 & bar \\
$T$ & Gas temperature & 300 & K \\
\bottomrule
\end{tabular}
\end{table}


% ══════════════════════════════════════════════════════════════════════════════
\section{Xenon Gas Properties}
\label{sec:gas}
% ══════════════════════════════════════════════════════════════════════════════

At \SI{15}{bar} and \SI{300}{K}:

\begin{table}[H]
\centering
\caption{Xenon gas properties.}
\label{tab:gas}
\begin{tabular}{lll}
\toprule
Property & Value & Note \\
\midrule
Density $\rho$ & \SI{79.0}{kg/m^3} & 66$\times$ denser than air \\
Dynamic viscosity $\mu$ & \SI{2.32e-5}{Pa\cdot s} & --- \\
Kinematic viscosity $\nu$ & \SI{2.94e-7}{m^2/s} & 51$\times$ smaller than air \\
Speed of sound $c_s$ & \SI{177.9}{m/s} & --- \\
\bottomrule
\end{tabular}
\end{table}


% ══════════════════════════════════════════════════════════════════════════════
\section{Three-Phase Motion Sequence}
\label{sec:phases}
% ══════════════════════════════════════════════════════════════════════════════

The complete motion consists of three sequential phases:

\begin{enumerate}[nosep]
  \item \textbf{Phase 1}: Blade rotation (bang-bang profile)
  \item \textbf{Phase 2}: Carrier radial slide (bang-bang profile)
  \item \textbf{Phase 3}: Vertical plate lift + wait period
\end{enumerate}


% ══════════════════════════════════════════════════════════════════════════════
\section{Phase 1: Blade Rotation}
\label{sec:phase1}
% ══════════════════════════════════════════════════════════════════════════════

\subsection{Kinematics}

The propeller executes a bang-bang motion profile under constant torque $\tau$:
\begin{align}
  I &= n \times \frac{1}{3}m_\mathrm{blade}R^2 + n \times m_\mathrm{plate}r_0^2 \\
  t_\mathrm{rot} &= 2\sqrt{\frac{\theta I}{\tau}} \\
  \omega_\mathrm{max} &= \sqrt{\frac{\theta\tau}{I}} \\
  V_\mathrm{tip} &= \omega_\mathrm{max} R
\end{align}

\begin{table}[H]
\centering
\caption{Phase 1 kinematics (2-arm propeller).}
\label{tab:phase1_kin}
\begin{tabular}{lll}
\toprule
Quantity & Value & Units \\
\midrule
Moment of inertia $I$ & 0.747 & kg$\cdot$m$^2$ \\
Angular acceleration $\alpha = \tau/I$ & 80.3 & rad/s$^2$ \\
Rotation time $t_\mathrm{rot}$ & 396 & ms \\
Peak angular velocity $\omega_\mathrm{max}$ & 15.9 & rad/s \\
Peak tip velocity $V_\mathrm{tip}$ & 25.4 & m/s \\
\bottomrule
\end{tabular}
\end{table}

\subsection{Flow Regime}

\begin{table}[H]
\centering
\caption{Flow regime at peak velocity.}
\label{tab:regime}
\begin{tabular}{lll}
\toprule
Quantity & Value & Interpretation \\
\midrule
Tip Mach number Ma & 0.143 & Incompressible ($< 0.3$) \\
Chord Reynolds number Re$_c$ & $1.38 \times 10^7$ & Fully turbulent ($> 5 \times 10^5$) \\
\bottomrule
\end{tabular}
\end{table}

\subsection{Boundary Layer}

For a turbulent BL developing over a flat plate of length $c$:
\begin{align}
  \delta &= 0.37\,c\,\mathrm{Re}_c^{-1/5} \\
  \delta^* &= \delta/8 \\
  C_f &= 0.0592\,\mathrm{Re}_c^{-1/5} \\
  u_\tau &= V_\mathrm{tip}\sqrt{C_f/2}
\end{align}

\begin{table}[H]
\centering
\caption{Blade boundary layer properties.}
\label{tab:BL}
\begin{tabular}{lll}
\toprule
Quantity & Value & Units \\
\midrule
BL thickness $\delta_1$ & 2.21 & mm \\
Displacement thickness $\delta^*$ & 0.28 & mm \\
Friction velocity $u_\tau$ & 0.84 & m/s \\
Turbulence intensity $u'$ & 1.27 & m/s \\
Eddy timescale $t_e = \delta/u'$ & 1.7 & ms \\
Decay time $2\,t_e$ & 3.5 & ms \\
\bottomrule
\end{tabular}
\end{table}


% ══════════════════════════════════════════════════════════════════════════════
\section{Phase 1 Mechanisms Analysis}
\label{sec:mechanisms}
% ══════════════════════════════════════════════════════════════════════════════

We analyze all potential sources of residual gas velocity from the blade rotation.

\subsection{Mechanism 1: Turbulent Eddy Diffusion}

After the blade stops, the momentum in the boundary layer can diffuse
vertically into the quiescent gas above. The evolution is governed by:
\begin{equation}
  \frac{\partial u}{\partial t} = \frac{\partial}{\partial z}
    \left[\nu_\mathrm{eff}(z,t)\frac{\partial u}{\partial z}\right]
\end{equation}

The turbulent viscosity decays exponentially after motion stops:
\begin{equation}
  \nu_t(z,t) = \nu_t(z,0)\exp\left(-\frac{t}{2t_e}\right)
\end{equation}

\begin{table}[H]
\centering
\caption{Turbulent diffusion parameters.}
\label{tab:diffusion}
\begin{tabular}{lll}
\toprule
Quantity & Value & Units \\
\midrule
Outer turbulent viscosity $\nu_t$ & $1.26 \times 10^{-4}$ & m$^2$/s \\
Ratio $\nu_t/\nu$ & 430 & --- \\
Turbulent reach $\ell_\mathrm{turb}$ & 0.66 & mm \\
Molecular reach $\ell_\mathrm{mol}$ & 0.38 & mm \\
Max reach from BL top & 3.3 & mm \\
\bottomrule
\end{tabular}
\end{table}

\textbf{Result:} Momentum is confined within $\sim$3 mm of the blade surface.


\subsection{Mechanism 2: Bulk Swirl and Angular Momentum}

The drag on the blades deposits tangential momentum into the gas:
\begin{equation}
  \tau_\mathrm{drag} = n \int_0^R r \cdot \tfrac{1}{2}\rho[\omega r]^2 c\, C_d\, dr
\end{equation}

\begin{table}[H]
\centering
\caption{Bulk swirl parameters.}
\label{tab:swirl}
\begin{tabular}{lll}
\toprule
Quantity & Value & Units \\
\midrule
Drag torque $\tau_\mathrm{drag}$ & 47.7 & N$\cdot$m \\
Angular impulse $J$ & 18.9 & N$\cdot$m$\cdot$s \\
Molecular diffusion time $t_\mathrm{visc}$ & 26.5 & s \\
Ekman spin-up time (molecular) & 736 & s \\
Ekman spin-up time (turbulent) & 39 & s \\
\bottomrule
\end{tabular}
\end{table}

\textbf{Result:} Momentum stays confined to blade plane. No transport
within the $\sim$2 s cycle time.


\subsection{Mechanism 3: Displacement (Potential) Flow}

The NACA\,0012 blade has finite thickness $t_\mathrm{max} = 0.12c$. As it
moves through the gas, it displaces fluid. When the blade stops, the
potential flow field propagates away at the speed of sound.

\begin{table}[H]
\centering
\caption{Potential flow parameters.}
\label{tab:potential}
\begin{tabular}{lll}
\toprule
Quantity & Value & Units \\
\midrule
Displacement velocity (during sweep) & 20.0 & m/s \\
Acoustic clearing time $t_\mathrm{acoustic}$ & 9.0 & ms \\
\bottomrule
\end{tabular}
\end{table}

\textbf{Result:} Potential flow clears in $\sim$9 ms. Zero residual.


\subsection{Mechanism 4: Pressure Pulse from Deceleration}

When the blade decelerates to rest, the change in momentum creates a
pressure transient:
\begin{align}
  \Delta p &\sim \rho V_\mathrm{tip}^2 \\
  \Delta u &\sim \frac{\Delta p}{\rho c_s} = \frac{V_\mathrm{tip}^2}{c_s}
\end{align}

\begin{table}[H]
\centering
\caption{Pressure pulse parameters.}
\label{tab:pressure}
\begin{tabular}{lll}
\toprule
Quantity & Value & Units \\
\midrule
Pressure perturbation $\Delta p$ & 51.0 & kPa \\
Fractional $\Delta p/P$ & 3.4\% & --- \\
Induced acoustic velocity $\Delta u$ & 3.6 & m/s \\
\bottomrule
\end{tabular}
\end{table}

\textbf{Result:} Transient pulse passes through observation region. Zero residual.


\subsection{Mechanism 5: Secondary Flows}

\paragraph{Tip Clearance Flow.}
For NACA\,0012 at $\alpha = 0$: there is \emph{no lift}, hence \emph{no
pressure difference} across the blade.

\paragraph{Centrifugal Effects.}
The tangential wake flows \emph{along} the cylindrical wall, not into it.
There is no face-on impingement to deflect flow into the axial direction.

\begin{table}[H]
\centering
\caption{Secondary flow parameters.}
\label{tab:secondary}
\begin{tabular}{lll}
\toprule
Quantity & Value & Units \\
\midrule
Tip gap & 5.0 & mm \\
Centrifugal acceleration & 101 & m/s$^2$ \\
\bottomrule
\end{tabular}
\end{table}

\textbf{Result:} No mechanism for axial momentum transport. Zero residual.


\subsection{Mechanism 6: Vortex Shedding}

For a \emph{non-lifting} blade ($\alpha = 0$), there is no net circulation,
hence no macroscopic starting/stopping vortex. NACA\,0012 at zero angle
of attack is streamlined: coherent vortex shedding is strongly suppressed.

\textbf{Result:} Any shed vortices confined to blade plane. Zero residual.


% ══════════════════════════════════════════════════════════════════════════════
\section{Phase 2: Carrier Radial Slide}
\label{sec:phase2}
% ══════════════════════════════════════════════════════════════════════════════

After rotation, the carrier plate slides radially along the blade arm.

\begin{table}[H]
\centering
\caption{Phase 2 kinematics.}
\label{tab:phase2_kin}
\begin{tabular}{lll}
\toprule
Quantity & Value & Units \\
\midrule
Slide distance $\Delta r$ & 0.8 & m \\
Slide time $t_\mathrm{slide}$ & 200 & ms \\
Peak velocity $V_\mathrm{slide}$ & 8.0 & m/s \\
\bottomrule
\end{tabular}
\end{table}

\begin{table}[H]
\centering
\caption{Carrier boundary layer properties.}
\label{tab:carrier_BL}
\begin{tabular}{lll}
\toprule
Quantity & Value & Units \\
\midrule
Carrier Reynolds number Re$_L$ & $4.4 \times 10^6$ & --- \\
BL thickness $\delta_2$ & 2.78 & mm \\
Decay time $2t_e$ & 14 & ms \\
Max reach & 4.0 & mm \\
\bottomrule
\end{tabular}
\end{table}

\textbf{Result:} Same analysis as blade BL. Momentum confined within
$\sim$4 mm of the carrier surface.


% ══════════════════════════════════════════════════════════════════════════════
\section{Phase 3: Vertical Lift and Quiescence}
\label{sec:phase3}
% ══════════════════════════════════════════════════════════════════════════════

\subsection{Lift Kinematics}

\begin{table}[H]
\centering
\caption{Phase 3 kinematics.}
\label{tab:phase3_kin}
\begin{tabular}{lll}
\toprule
Quantity & Value & Units \\
\midrule
Lift height $z_\mathrm{lift}$ & 5.0 & mm \\
Lift duration $t_\mathrm{lift}$ & 1.0 & s \\
Lift acceleration $a_\mathrm{lift}$ & 0.020 & m/s$^2$ \\
Peak lift velocity $V_\mathrm{lift}$ & 10.0 & mm/s \\
Wait time $t_\mathrm{wait}$ & 0.5 & s \\
Stokes layer thickness $\delta_3$ & 0.54 & mm \\
\bottomrule
\end{tabular}
\end{table}

\subsection{Three Velocity Profiles}

After the lift, we evaluate three velocity profiles to determine quiescence:

\paragraph{Profile 1: Blade BL Diffusion.}
The blade boundary layer from Phase 1, evaluated at time
$t = t_\mathrm{slide} + t_\mathrm{lift} + t_\mathrm{wait} = \SI{1.7}{s}$
after rotation stopped. Since the turbulence decay time is only \SI{3.5}{ms},
this profile is \textbf{fully decayed} by the evaluation time.

\paragraph{Profile 2: Carrier Slide BL Diffusion.}
The carrier boundary layer from Phase 2, evaluated at time
$t = t_\mathrm{lift} + t_\mathrm{wait} = \SI{1.5}{s}$ after the slide stopped.
With a decay time of \SI{14}{ms}, this profile is also essentially decayed.

\paragraph{Profile 3: Stokes Layer from Lift.}
The boundary layer created by the vertical lift, evaluated at $t_\mathrm{wait}$
after the lift stops. This follows the classical Stokes first problem
(impulsively stopped plate):
\begin{equation}
  u(z, t) = V_\mathrm{lift}\,\mathrm{erfc}\!\left(\frac{z - z_\mathrm{lift}}{2\sqrt{\nu_\mathrm{eff}\,t}}\right)
\end{equation}

\subsection{Quiescence Heights}

\begin{table}[H]
\centering
\caption{Quiescence heights for 2-arm propeller ($u_c = \SI{0.1}{mm/s}$).}
\label{tab:quiescence}
\begin{tabular}{llll}
\toprule
Profile & Description & $z_q$ & Note \\
\midrule
Profile 1 & Blade BL & \SI{0.1}{mm} & Fully decayed \\
Profile 2 & Carrier slide BL & \SI{2.8}{mm} & Essentially decayed \\
Profile 3 & Stokes layer & \SI{7.5}{mm} & Dominant contribution \\
\bottomrule
\end{tabular}
\end{table}

The dominant contribution is Profile 3 (Stokes layer from the vertical lift).
Since the plate surface is at $z = z_\mathrm{lift} = \SI{5}{mm}$, the
quiescence height $z_{q3} = \SI{7.5}{mm}$ is only \textbf{\SI{2.5}{mm} above
the plate surface}.


% ══════════════════════════════════════════════════════════════════════════════
\section{Summary of All Mechanisms}
\label{sec:summary_mechanisms}
% ══════════════════════════════════════════════════════════════════════════════

\begin{table}[H]
\centering
\caption{Summary of all mechanisms for residual velocity (2-arm propeller).}
\label{tab:all_mechanisms}
\begin{tabular}{llc}
\toprule
Mechanism & Timescale & Result \\
\midrule
\multicolumn{3}{l}{\textit{Phase 1: Blade rotation}} \\
1. Eddy diffusion & $2t_e \approx 3.5$ ms & $3.2 \times 10^{-7}$ m/s \\
2. Bulk swirl & $t_\mathrm{visc} \approx 26$ s & 0 \\
3. Potential flow & $t_\mathrm{acoustic} \approx 9$ ms & 0 \\
4. Pressure pulse & $t_\mathrm{decel} \sim 10$ ms & 0 \\
5. Secondary flows & $t_\mathrm{Ekman} \approx 39$ s & 0 \\
6. Vortex shedding & N/A (confined) & 0 \\
\midrule
\multicolumn{3}{l}{\textit{Phase 2: Carrier slide}} \\
7. Carrier eddy diff. & max reach $<$ 5 mm & 0 \\
\midrule
\multicolumn{3}{l}{\textit{Phase 3: Vertical lift (three profiles)}} \\
8. Blade BL ($z_{q1}$) & eval at $t = 1.7$ s & \SI{0.1}{mm} \\
9. Carrier BL ($z_{q2}$) & eval at $t = 1.5$ s & \SI{2.8}{mm} \\
10. Stokes layer ($z_{q3}$) & eval at $t_\mathrm{wait}$ & \SI{7.5}{mm} \\
\bottomrule
\end{tabular}
\end{table}


% ══════════════════════════════════════════════════════════════════════════════
\section{Timeline and Conclusion}
\label{sec:conclusion}
% ══════════════════════════════════════════════════════════════════════════════

\subsection{Timeline}

\begin{table}[H]
\centering
\caption{Complete motion timeline for 2-arm propeller.}
\label{tab:timeline}
\begin{tabular}{lll}
\toprule
Phase & Duration & Cumulative \\
\midrule
Phase 1: Blade rotation & \SI{396}{ms} & \SI{0.40}{s} \\
Phase 2: Carrier slide & \SI{200}{ms} & \SI{0.60}{s} \\
Phase 3: Vertical lift & \SI{1000}{ms} & \SI{1.60}{s} \\
Wait period & \SI{500}{ms} & \SI{2.10}{s} \\
\midrule
\textbf{Total cycle time} & & \textbf{\SI{2.1}{s}} \\
\bottomrule
\end{tabular}
\end{table}

\subsection{Key Results}

\begin{table}[H]
\centering
\caption{Summary of key parameters for 2-arm propeller.}
\label{tab:summary}
\begin{tabular}{lll}
\toprule
Parameter & Value & Units \\
\midrule
\multicolumn{3}{l}{\textit{Peak velocities}} \\
Blade tip velocity $V_\mathrm{tip}$ & 25.4 & m/s \\
Carrier slide velocity $V_\mathrm{slide}$ & 8.0 & m/s \\
Lift velocity $V_\mathrm{lift}$ & 10.0 & mm/s \\
\midrule
\multicolumn{3}{l}{\textit{Boundary layer thicknesses}} \\
Blade BL $\delta_1$ & 2.21 & mm \\
Carrier BL $\delta_2$ & 2.78 & mm \\
Stokes layer $\delta_3$ & 0.54 & mm \\
\midrule
\multicolumn{3}{l}{\textit{Quiescence heights}} \\
$z_{q1}$ (blade BL) & 0.1 & mm \\
$z_{q2}$ (carrier BL) & 2.8 & mm \\
$z_{q3}$ (Stokes layer) & 7.5 & mm \\
\midrule
\multicolumn{3}{l}{\textit{Timing}} \\
Total cycle time & 2.1 & s \\
\bottomrule
\end{tabular}
\end{table}

\subsection{Conclusion}

For the 2-arm propeller configuration, the gas above the lifted carrier
plate becomes \textbf{quiescent} (velocity $< \SI{0.1}{mm/s}$) within the
total cycle time of \textbf{\SI{2.1}{s}}.

The fundamental reasons for rapid quiescence are:
\begin{enumerate}[nosep]
  \item The NACA\,0012 at zero angle of attack produces \textbf{no lift},
    so all aerodynamic effects are confined to the thin drag wake within
    $\delta \approx 2$--3 mm of the surface.
  \item Dense xenon ($\rho \approx \SI{79}{kg/m^3}$) has very small
    kinematic viscosity ($\nu \approx \SI{3e-7}{m^2/s}$), making diffusion
    timescales very long---momentum stays confined.
  \item Turbulent eddies decay rapidly ($\sim$\SI{3.5}{ms}) once motion stops.
    By evaluation time ($\sim$1.5--1.7 s), both blade and carrier boundary
    layers have fully decayed.
  \item The slow lift velocity ($V_\mathrm{lift} = \SI{10}{mm/s}$) creates
    only a thin Stokes layer ($\delta_3 \approx \SI{0.5}{mm}$).
  \item The dominant residual flow is Profile 3 (Stokes layer), which gives
    quiescence at $z_{q3} = \SI{7.5}{mm}$, just \textbf{\SI{2.5}{mm} above
    the lifted plate surface}.
\end{enumerate}

\end{document}
